% 概念名词表为整页长表：若本节标题落在页底，表体只会有两三行溢出到次页，
% 既难看又多占一页；故整表另页起排（同口径 19：长表另页起排的处置）。
\clearpage
\section{概念名词表}

本文用到的游戏概念名词见表 \ref{tab:conceptnouns}。

{\small
\begin{longtable}{@{}>{\raggedright\arraybackslash}p{4.7cm}>{\raggedright\arraybackslash}p{2.4cm}>{\raggedright\arraybackslash}p{8.5cm}@{}}
\caption{概念名词表}\label{tab:conceptnouns}\\
\toprule
\textbf{名词} & \textbf{归属} & \textbf{释义} \\
\midrule
\endfirsthead
\toprule
\textbf{名词} & \textbf{归属} & \textbf{释义} \\
\midrule
\endhead
\bottomrule
\endlastfoot
\term{月兆}（Moonsign） & 月曜体系 & 队伍加成系统；由月兆角色数量决定等级（初辉／满辉） \\
\rowline
\term{月兆祝赐}（Moonsign Benediction） & 月曜体系 & “月兆”相关固有天赋的命名前缀（如“月兆祝赐·千籁恩宠”）；共同效果为“在队伍中时，队伍的月兆将会上升一级” \\
\rowline
\term{草露}（Verdant Dew） & \rxLunar{月绽放} & 触发\rxLunar{月绽放}后 2.5 秒填充 1 枚的队伍资源，供特定角色消耗 \\
\rowline
\term{诳言之核}（Seeds of Deceit） & \rxLunar{月绽放} & 奈芙尔把草原核／丰穰之核转化成的资源核，不自爆、可被其重击吸收 \\
\rowline
\term{伪秘之帷}（Veil of Falsehood） & \rxLunar{月绽放} & 奈芙尔吸收诳言之核获得的层数增益（至多 3 层；元素爆发消耗层数、每层增伤 13\%→49\%；满 3 层时元素精通 +100，持续 8 秒） \\
\rowline
\term{雷暴云}（Thundercloud） & \rxLunar{月感电} & \rxLunar{月感电}反应的载体，每 2 秒对同时附有水雷的目标雷击 \\
\rowline
\term{月笼}（Moondrift） & \rxLunar{月结晶} & \rxLunar{月结晶}反应生成的三枚岩元素造物 \\
\rowline
\term{月笼谐奏}（Moondrift Harmony） & \rxLunar{月结晶} & 每 3 次\rxLunar{月结晶}引发的范围岩元素伤害 \\
\rowline
\term{星之楔}（Stellar Linchpin） & 星烁体系 & 部分星烁体系角色的属性栏条目（取代“神之眼”一栏；7.0 及之前可玩角色为桑多涅、奥黛塔） \\
\rowline
\term{星辉棱晶}（Starlight Prism） & \rxStellar{星超导} & 触发\rxStellar{星超导}时生成、并以此为心展开极星辉域的结构 \\
\rowline
\term{极星辉域}（Polestar Field） & \rxStellar{星超导} & 持续 6 秒的领域，提供冰/雷增伤、物抗降低与\rxStellar{星超导}系数提升（最高 2.0） \\
\rowline
\term{星辉风旋}（Stellar Vortex） & \rxStellar{星扩散} & \rxStellar{星扩散}生成的风旋，至多 1 枚，3 秒后（或第 6 次触发时提前）引爆 \\
\rowline
\term{辉映}（Radiance） & 星烁体系 & 「辉映·\rxStellar{星超导}」「辉映·\rxStellar{星扩散}」，强化特定角色技能的状态；\rxStellar{星超导}优先 \\
\rowline
\term{擢升}（Elevation） & 月曜／星烁 & 独立最终乘区，目前仅由限定五星命之座提供 \\
\rowline
\term{魔女的谕示}（Witch’s Revelation） & 星烁体系 & 月之八（6.7）起赋予部分老角色的强化系统：在「\term{极星辉域}」内获得追加强化，部分角色可打出星烁反应伤害，但其本身不能开启星烁反应 \\
\rowline
\term{直伤型 / 反应型}（Direct / Reaction） & 月曜·星烁 & 直伤型 = 角色天赋直接打出、被“视为”该反应伤害；反应型 = 由反应衍生物（雷暴云／月笼谐奏／风旋爆炸）结算。（英文为社区分类词，官方文本只有“视为该反应伤害”的句式表述） \\
\rowline
\term{激元素}（Quicken） & 激化 & 原激化施加的状态元素（约 6～9.5 秒）；官方英文只有 Quicken／Quickened，没有“元素”这一名词层级 \\
\rowline
\term{丰穰之核}（Bountiful Core） & 绽放 & 妮露专属的强化草原核 \\
\end{longtable}
}

\medskip

\clearpage

\section{参考资料与信源分级}

\subsection{分级标准}\label{sec:srcgrade}

信源分为 T0–T4 五级，定义与本文用法见表 \ref{tab:srcgrade}。

\begin{table}[ht]\centering\small
\caption{信源分级标准}\label{tab:srcgrade}
\begin{tabular}{@{}l>{\raggedright\arraybackslash}p{7cm}>{\raggedright\arraybackslash}p{7cm}@{}}
\toprule
\textbf{等级} & \textbf{定义} & \textbf{本文用法} \\
\midrule
\textbf{T0} & 官方文本：版本更新说明、官方博客、游戏内说明 & 机制定义的唯一依据；可直接引用 \\
\rowline
\textbf{T1} & 专业理论站（有成体系的机制文档、被广泛引用）：KeqingMains 等 & 公式结构与数值表；与 T2 交叉验证后写入 \\
\rowline
\textbf{T2} & 带\textbf{实测数据／拟合举证}的社群研究：逐乘区拟合、Excel 复算、样本统计 & 数值结论的主要来源；须说明误差与样本 \\
\rowline
\textbf{T3} & 社群 wiki 与媒体整理（无独立实测，但长期维护、可追责） & 仅作旁证与补齐描述 \\
\rowline
\textbf{T4} & 解包／爆料搬运（beta 数据、Seele\_Leaks 等）：信息有时效性，可能与上线版本不同 & 仅作时间线与旁证，不作数值依据 \\
\rowline
\textbf{不可用} & 注水内容：无数据、无口径、纯搬运的“爆料帖/标题党视频/AI 洗稿” & 一律不采信 \\
\bottomrule
\end{tabular}
\end{table}

\indent\textbf{三条件并列原则}（判定某位作者的新反应理论是否可靠）：① 有大量实测数据或实战细节描述；② 有（或曾有、且被广泛认可）一套理论体系；③ 有相当多人认可。三者缺一即降级。

\subsection{本次采信的信源}

本文实际采信的信源清单见表 \ref{tab:sources}。

{\small
\begin{longtable}{@{}>{\raggedright\arraybackslash}p{1.9cm}>{\raggedright\arraybackslash}p{6.5cm}>{\raggedright\arraybackslash}p{8.5cm}@{}}
\caption{本次采信的信源}\label{tab:sources}\\
\toprule
\textbf{等级} & \textbf{信源} & \textbf{用到的内容} \\
\midrule
\endfirsthead
\toprule
\textbf{等级} & \textbf{信源} & \textbf{用到的内容} \\
\midrule
\endhead
\bottomrule
\endlastfoot
T0 & 官方更新说明：3.0、5.8、月之一、月之四、月之八、7.0（简中官网 ys.mihoyo.com／繁中与英文官网 genshin.hoyoverse.com） & 五个新反应的官方定义、月兆系统原文、历史调整；概念名词表的英文名亦取自英文版官方文本（更新说明＋游戏内术语词典） \\
\rowline
T0 & 官方更新说明（月之四，哥伦比娅条目） & 哥伦比娅可使队伍同时触发\rxLunar{月感电}／\rxLunar{月绽放}／\rxLunar{月结晶}（三系月曜反应的共同开启者） \\
\rowline
T1 & KeqingMains《Lunar Reaction Guide》《Stellar Glimmer Reaction Guide》《Snezhnaya 版本综述》 & \rxLunar{月感电}/\rxLunar{月结晶}/\rxStellar{星超导}/\rxStellar{星扩散}的公式结构、\rxStellar{星超导}层数表、翠绿之影与深廊终曲等装备联动、月兆上限 \\
\rowline
T2 & bwiki 角色页（菈乌玛 / 奈芙尔）与 bwiki 攻略页 & 技能表倍率（284.5\%／每枚草露 273.6\%）、天赋原文（初辉 15\%／100\%、满辉 +10\%／+20\%、伪秘之帷 13\%→49\%、诳言之核「无法触发超绽放与烈绽放、也无法迸发」） \\
\rowline
T2 & 17173《菈乌玛》攻略（米游社转载） & 长按战技「消耗 3 枚草露 → 284.5\% 草伤 + 820.8\% \rxLunar{月绽放}伤害」、「霜林圣域」 \\
\rowline
T3 & Fandom 版本表（中文 wiki） & 「月之N ＝ 6.(N-1)」编号映射与各版本日期（\textbf{wiki 口径}；官方文本不带 6.x 编号） \\
\rowline
T1 & meropide（梅信心）《星烁反应》《伤害公式》中文理论速查站 & \rxStellar{星超导}叠层规则（一次攻击多目标只记 1 层、双判算 2 层、Buff 消失重置不继承）、风旋等级系数模型（1～6、达 6 立即引爆、最晚 3 秒）、高跳 6 秒、对群按目标累加 X 层、各乘区公式；留档 \texttt{调研/raw/补\_meropide星烁反应参考.md}、\texttt{补\_meropide伤害公式速查.md} \\
\rowline
T1 & KeqingMains 机制库《Snapshot \& Dynamic》《Transformative Reactions》(TCL) & 「剧变反应自 1.4 起不快照精通、实时更新」及奥兹／超载实例；留档 \texttt{调研/raw/补\_KQM快照机制\_snapshot-and-dynamic.md}、\texttt{补\_KQM-TCL剧变反应\_transformative-reactions.md} \\
\rowline
T1 & KeqingMains《Lunar Reaction Guide》公式图（本文抓取原图并逐字视觉转录） & 单角色公式（\rxLunar{月感电} $3\times$、\rxLunar{月结晶} $1.6\times$ 等级系数）、最终伤害权重 0.6／0.3／0.05／0.05 \\
\rowline
T2 & 幕陵《〈原神〉\rxLunar{月感电}机制解析》（游民星空转载，含 10 点拟合与三人实测样例） & \rxLunar{月感电}乘区、$6\text{EM}/(\text{EM}+2000)$ 拟合、1/1/2/1/12 加权、雷暴云 6 秒与 0.8U 支撑 2 次进程 \\
\rowline
T2 & 巴哈姆特实测帖（直伤\rxLunar{月结晶}公式 Excel 复算，误差 < 0.02\%） & \rxLunar{月结晶} 1.6 系数、精通与月增伤同乘区、羽毛不吃增伤 \\
\rowline
T2 & 4399《奥黛塔全面解析》（含\rxStellar{星超导}逐层实测表、\rxStellar{星扩散}触发统计） & \rxStellar{星超导}层数-系数-增伤对应表、\rxStellar{星扩散} 0.75/2/3 倍率与 0.6/0.3/0.05 分配 \\
\rowline
T2 & 电脑王阿达 kocpc（\rxStellar{星超导}乘区归纳） & \rxStellar{星超导}不吃防御、可暴击、擢升乘区 \\
\rowline
T2 & NGA 实测讨论帖《梦见月瑞希配队疑问…\rxStellar{星扩散}能否吃到珐露珊减风抗》(tid=47366657) & \rxStellar{星扩散}风伤段吃风减抗与元素暴伤、不吃元素伤害加成；风绿／冰蓝数字对应两段元素 \\
\rowline
T4（解包搬运，仅作旁证） & NGA 搬运帖 tid=47086461（风套调整/\rxStellar{星扩散}）、tid=44729195（\rxLunar{月绽放}数字与自伤） & 风套早期文本；\rxLunar{月绽放}艺术字与自伤的来源区分 \\
\rowline
\textbf{T3}（数值另经 T1/T2 交叉佐证） & bwiki《元素反应》《翠绿之影》 & 倍率总览、反应强度表、风套现行文本 \\
\rowline
T3 & 17173 / TapTap / 米游社攻略、9game、值得买转载 & 草露机制、菈乌玛与奈芙尔的直伤公式与数值、月兆双暴加成 \\
\end{longtable}
}

\subsection{取证受阻}

取证过程受阻的站点与处理方式见表 \ref{tab:blocked}。

\begin{table}[ht]\centering\small
\caption{取证受阻记录}\label{tab:blocked}
\begin{tabular}{@{}>{\raggedright\arraybackslash}p{3.4cm}>{\raggedright\arraybackslash}p{5.8cm}>{\raggedright\arraybackslash}p{6.4cm}@{}}
\toprule
\textbf{站点} & \textbf{状态} & \textbf{处理} \\
\midrule
NGA（bbs.nga.cn） & 无登录时返回 403「访客不能直接访问」；\textbf{已由系统 Edge（CDP 接入）登录后取用} & 已采信 3 帖并留档（\texttt{调研/raw/nga\_} 前缀的 3 个 txt）：tid=47366657（\rxStellar{星扩散}两段伤害与减抗，玩家实测）、tid=47086461（风套/\rxStellar{星扩散}，解包搬运）、tid=44729195（\rxLunar{月绽放}数字与自伤，解包搬运） \\
\rowline
Fandom 英文 wiki & 网络层不可达（\texttt{?action=raw} 亦仅回 862 字符） & 跳过 \\
\rowline
米游社文章页 & 前端渲染，正文不随 HTML 返回 & 以 17173/TapTap/4399 等转载镜像替代 \\
\bottomrule
\end{tabular}
\end{table}

\srcfoot{［留档完整性说明］所有引用一律以 \texttt{调研/raw/} 内文件为准；凡该目录无对应文字者一律列入 \sref{sec:open} 并注明状态，不混作已证。}

\subsection{存疑清单}\label{sec:open}

\begin{enumerate}
\item \textbf{月曜／星烁反应的精确元素量消耗}：\rxLunar{月感电}已确认为「每次雷击从水、雷两条元素量上\textbf{各扣 0.4U}」（等价于沿用常规感电的 0.4U 口径，但官方更新说明未给元素量数据，仅有 T1 口径）；\rxLunar{月结晶}「多段触发」的元素量细节仍无一致说法。\src{T1 KeqingMains 月曜指南与元素计量学｜T2 幕陵\rxLunar{月感电}实测}
\item \textbf{\rxLunar{月结晶}反应伤害是否参与防御区}：［T1］KeqingMains 称月曜反应不吃防御；部分中文整理文将其描述为“按防御力计算”，本文按前者（不吃防御）处理，后者可能与\textbf{直伤}\rxLunar{月结晶}（兹白/哥伦比娅的防御倍率技能）混淆。
\item \textbf{伪秘之帷自身的持续时间}：bwiki 角色页与攻略页均未标注；网络联查得 T3 多源一致为「9 秒、每层独立计算」（引游戏文本；出处记录见 调研/raw/补-伪秘之帷时长-web检索.md），未获 T0/T2。
\item \textbf{满辉与初辉的共处关系（表述分歧）}：官方仅表述「不同月兆等级提供的强化效果无法叠加」；bwiki 称满辉时也会触发初辉的对应效果，另有社区整理文称二者不可共存（该说法未留档）——本文将官方表述作为结论，不采纳任一方补充解读。
\item \textbf{\term{极星辉域}与风旋高跳的时长}：\term{极星辉域}本文取 \textbf{6 秒}（［T2］4399、［T3］bwiki／Mobalytics），［T1］meropide 实测作「约 7 秒」（注明 60 帧环境）；风旋爆炸的高跳本文取 \textbf{6 秒}（［T1］meropide、［T2］4399），［T1］KeqingMains 作 5 秒。官方对两者均只作「短暂」表述。
\item \textbf{\term{月兆·初辉}与\term{月兆·满辉}的等级口径}：官方原文只写「编入 1/2 名月兆角色时，分别使队伍的月兆等级变为初辉/满辉」，本文概念表按攻略口径写作「2 名及以上」。
\item \textbf{月曜反应对 90 级以下角色的等级系数}：本文仅记录 90 级（1446.85）口径。
\item \textbf{风套 4 件套的\rxStellar{星扩散}数值}：现行 wiki／数据站文本为“扩散 +60\%、\rxStellar{星扩散} +20\%、触发\rxStellar{星扩散}降 40\% 冰抗"，而早期搬运帖为"扩散、\rxStellar{星扩散} +60\%"；两版并存，本文按现行文本记录。
\item \textbf{冻结的\textbf{冻结元素量}算法}：［T1］KeqingMains 记“$2\times$ 双方元素量的较小者”，［T3］bwiki 记“反应元素量的\textbf{总和}”；两说互不相同，本文仅记录该分歧，不在消耗比表内给冻结的比值。
\end{enumerate}

\subsection{取值与留档说明}\label{sec:notes}

\begin{itemize}
\item \textbf{官方数值优先}：能取到官方原文的一律以官方为准（版本日期、反应定义、属性栏条目、装备效果数值）。
\item \textbf{社区数值的来源级别}：月曜／星烁的系数、层数—增益对应表、加权比例等，均出自 \src{T1}（KeqingMains、meropide）与 \src{T2}（带实测举证）两级；单一来源支持的结论在正文逐条标注出处，孤例与多源冲突者一律列入上一节。
\item \textbf{留档}：所有引用一律以 \texttt{调研/raw/} 内文件为准；凡该目录无对应文字者一律列入上一节并注明状态，不混作已证。
\item \textbf{官方文本不带 6.x 编号}：正文中的「月之N = 6.(N-1)」对应关系出自 \textbf{Fandom 版本表}（wiki 口径），仅作注释，不视为官方结论（官方文本自称「月之N」／Luna I–VIII；6.1～6.8 在官方全量公告中零出现，仅“6.0 版本”作为对下一大版本的前向指称出现过）。
\item \textbf{“0.96”的出处}：中文来源常引用的“反应\rxLunar{月结晶}等效系数 0.96 = 1.6 × 0.6”的原始出处（早期检索时记为某篇 IGN 中国攻略）未能复现；本文采用该值，依据是与 KeqingMains 公式图、贡献权重 0.6 及\rxLunar{月感电} 1.8 的独立推导一致。
\item \textbf{未收录的社群说法}：“感电不再打断攻击”在社群 wiki 的历史记录里被列为“未知版本”，但官方公告归档（1193 篇，2019-06～2026-09）中没有“感电”与“打断”并列的条目，版本与出处无法落实，本文不予收录（官方唯一同类条目是 1.2 对「阴燃之火」持续伤害打断效果的移除）。
\end{itemize}

